\tikzset{add/.style n args={4}{
    minimum width=14,
    path picture={
    \draw[black] 
        (path picture bounding box.south east) -- (path picture bounding box.north west)
        (path picture bounding box.south west) -- (path picture bounding box.north east);
    \node[anchor=north, inner sep=0pt] at ($(path picture bounding box.north)+(0,-0.13)$) {\tiny #1};
    \node[anchor=east,  inner sep=0pt] at ($(path picture bounding box.east)+(-0.13,0)$)  {\tiny #2};
    \node[anchor=south, inner sep=0pt] at ($(path picture bounding box.south)+(0,0.13)$) {\tiny #3};
    \node[anchor=west,  inner sep=0pt] at ($(path picture bounding box.west)+(0.13,0)$)  {\tiny #4};
        }
    }
}

\tikzstyle{block} = [draw, rectangle, align=center]%
\tikzstyle{sum3} = [draw, circle, add={$+$}{$+$}{$+$}{}]%
\tikzstyle{sum2} = [draw, circle, add={}{$+$}{$+$}{}]%
\tikzstyle{comp} = [draw, circle, add={$-$}{$+$}{}{}]%
\tikzstyle{input} = [text centered, inner xsep=0pt, outer sep=0pt]%
\tikzstyle{output} = [text centered, inner xsep=0pt, outer sep=0pt]%

\begin{tikzpicture}[auto,node distance=1cm,>=latex']
    % ---------------------------------
	% --- taille de font par défaut ---
	% ---------------------------------

	% ----------------------------------
	% --- définition des espacements ---
	% ----------------------------------
	\def\x{1.5em}
    \def\y{1em}

    % ----------------------------
	% --- définition des nodes ---
	% ----------------------------
	% node de consigne
	\node [input] (r2prime) {$r'_2(k)$};

    % node de l'interpolateur
    \node [block, right=\x of r2prime.east] (interpo) {Interpolateur};

    % node du comparateur de consigne
    \node [comp, right=2*\x of interpo.east] (compR2) {};

    % node de l'intégrateur
    \node [block, right=\x of compR2.east] (Cp) {$C_P(q)$};

    % node du proportionnel
    \node [block, below=\y of Cp.south] (Ci) {$C_i(q)$};

    % node de l'additionneur
    \node [sum3, right=\x of Cp.east] (add) {};

    % node du procédé
    \node [block, right=\x of add.east] (G) {$G_o(q)$};

    % node de l'additionneur
    \node [sum2, right=\x of G.east] (addH) {};

    % node du modèle de bruit
    \node [block, above=\y of addH.north] (H) {$H_o(k)$};

    % node du bruit blanc
    \node [input, above=\y of H.north] (e) {$e_o(k)$};
    
    % node de la sortie
    \node [output, right=\x of addH.east] (y) {$y(k)$};

    % node du dérivateur
    \node [block, above=\y of compR2.north] (deriv) {$\frac{\text{d}}{\text{d}t}$};

    % node de l'inertie
    \node [block,  right=\x of deriv.east] (J) {$J$};

    % ------------------------------
	% --- définition des flèches ---
	% ------------------------------
    % de r'2 à l'interpolateur
    \draw [->] (r2prime) -- (interpo);

    % de l'interpolateur au comparateur
    \draw [->] (interpo) -- (compR2) node [midway, below] (r2) {$r_2(k)$};

    % de r2 au dérivateur
    \draw [->] (r2) |- (deriv);
    
    % du dérivateur à l'inertie
    \draw [->] (deriv) -- (J);

    % de l'inertie J à l'additionneur
    \draw [->] (J) -| (add) node [midway] {$r_1(k)$};

    % du comparateur aux régulateurs
    \draw [->] (compR2) -- (Cp) node [midway] (epsilon) {};
    \draw [->] (epsilon) |- (Ci);

    % des régulateurs à l'additionneur
    \draw [->] (Cp) -- (add);
    \draw [->] (Ci) -| (add);

    % de l'additionneur au procédé
    \draw [->] (add) -- (G);

    % du procédé à l'additionneur H
    \draw [->] (G) -- (addH);

    % du bruit au modèle de bruit
    \draw [->] (e) -- (H);

    % du modèle de bruit à l'additionneur H
    \draw [->] (H) -- (addH);

    % de l'additionneur H à la sortie
    \draw [->] (addH) -- (y) node [midway] (y2) {};

    % de l'additionneur au comparateur
    \draw [->] (y2) |- ([yshift=-\y]Ci.south) -| (compR2);
    
\end{tikzpicture}